% En el preámbulo: \usepackage{booktabs,longtable,array}
\begingroup
\small
\renewcommand{\arraystretch}{1.2}
\setlength{\tabcolsep}{4pt}
\begin{longtable}{@{}
  >{\raggedright\arraybackslash}p{0.21\linewidth}
  >{\raggedright\arraybackslash}p{0.15\linewidth}
  >{\raggedright\arraybackslash}p{0.23\linewidth}
  >{\raggedright\arraybackslash}p{0.33\linewidth}@{}}
\caption{Stack tecnológico y justificación de selección}
\label{tab:stack-tecnologico}\\
\toprule
\textbf{Componente} & \textbf{Versión} & \textbf{Función} & \textbf{Criterio de selección} \\
\midrule
\endfirsthead
\multicolumn{4}{c}{\tablename\ \thetable\ (continuación)}\\
\toprule
\textbf{Componente} & \textbf{Versión} & \textbf{Función} & \textbf{Criterio de selección} \\
\midrule
\endhead
\midrule
\multicolumn{4}{r}{Continúa en la siguiente página.}\\
\endfoot
\bottomrule
\endlastfoot
Angular & 21.2.23 &
Interfaz web de gestión clínica. &
Permite organizar formularios, navegación y vistas mediante componentes reutilizables para la administración de pacientes y solicitudes radiográficas. \\
TypeScript & 5.9.3 &
Lenguaje de implementación del frontend web. &
El tipado estático facilita detectar inconsistencias entre formularios, modelos y respuestas de la API durante el desarrollo. \\
Angular Material & 21.2.14 &
Componentes visuales del panel web. &
Proporciona controles uniformes para formularios, tablas, diálogos y navegación, reduciendo el trabajo de construcción de la interfaz. \\
Node.js / Express &
24.16.0 / 5.2.1 &
API HTTP, validaciones y lógica de negocio. &
Centraliza la autenticación, los permisos y el acceso a los datos. Su API proporciona el punto de integración definido para los componentes web y móvil con las mismas reglas de autorización. \\
PostgreSQL en Supabase & 17.6 &
Persistencia de datos clínicos y administrativos. &
El modelo relacional permite representar pacientes, odontólogos, asignaciones y solicitudes mediante claves, restricciones y transacciones que conservan la consistencia de los registros. \\
Supabase Storage &
Servicio administrado; SDK JavaScript 2.117.2 &
Almacenamiento de archivos radiográficos. &
Permite almacenar los archivos en un bucket privado y utilizar URL firmadas para su carga y consulta, separando los archivos de sus metadatos relacionales. \\
jsonwebtoken / bcryptjs & 9.0.3 / 3.0.3 &
Gestión de sesiones y verificación de contraseñas. &
JWT identifica las sesiones en las solicitudes a la API; bcrypt permite verificar contraseñas mediante hashes. Los permisos y las asignaciones se comprueban en el backend. \\
Flutter &
SDK por verificar en el proyecto móvil &
Interfaz de la APK Android complementaria. &
Se selecciona para adaptar la aplicación móvil existente a un alcance acotado: inicio de sesión del odontólogo, consulta de pacientes asignados y visualización de sus radiografías. \\
Dart &
SDK por verificar junto con Flutter &
Lenguaje de implementación de la aplicación móvil. &
Su integración con Flutter permite implementar la interfaz y las operaciones asíncronas necesarias para consumir la API y consultar los archivos autorizados. \\
Vercel &
Plataforma administrada &
Despliegue de la aplicación web y su API. &
La configuración del proyecto permite publicar el frontend y el backend y dirigir las solicitudes de la interfaz hacia el servicio de la API. \\
Git / GitHub &
Git 2.44.0.windows.1; GitHub como servicio &
Control de versiones y alojamiento del repositorio. &
Conserva el historial de cambios del código, las migraciones y la documentación, facilitando la recuperación de versiones y la revisión del avance del proyecto. \\
\end{longtable}
\noindent\textit{Fuente: Elaboración propia a partir del código, el archivo de dependencias bloqueadas, el entorno local y los metadatos de la base de datos.}

\medskip
\noindent\textbf{Alcance del componente móvil.}
La APK tiene como alcance la consulta para el odontólogo:
inicio y cierre de sesión, consulta de pacientes asignados, consulta del
estado de solicitudes y visualización de radiografías disponibles.
La administración de cuentas, roles y permisos y la gestión de odontólogos,
horarios y citas permanecerán en la aplicación web.
El paciente será un registro clínico y no iniciará sesión en la APK;
CERPAX continuará utilizando el enlace temporal de entrega.

\medskip
\noindent\textbf{Integración y versiones.}
El diseño de integración del componente móvil establece la API de la
plataforma como punto de acceso a los datos y de control de autorización.
La entrega móvil revisada se identifica como DentalCare 1.0.0+1, dato que
corresponde a la aplicación y no a las versiones de Flutter o Dart.
Estas dos versiones deberán confirmarse en el proyecto móvil antes de
cerrar la monografía. La versión de Node.js indicada corresponde al entorno
local verificado, no a una comprobación del entorno de producción.
\endgroup
